%% ==========================================================================
%%  Approach 7 --- local search on a lexicographic potential.
%%
%%  Source: derived 2026-08-06.  Scripts: update_18/localsearch_lemma.py,
%%  potentials.py, stress_potentials.py, p6_deep.py.
%%
%%  STATUS (2026-08-06).  PROVED: the reduction --- the descent lemma (L_P6)
%%  implies BOTH halves of Conjecture conj:main, existence and polynomial
%%  computability, with an explicit O(m^4 n^4) bound.  PROVED: the sub-case
%%  split, which isolates balanced partitions as the entire remaining burden.
%%  OPEN: (L_P6) itself, backed by ZERO stuck partitions over every partition
%%  of every instance tested (n <= 5, m <= 9, including 253 EF-free instances
%%  and R10's certified-hard family), and 7,416 crux partitions.
%%  REFUTED: the cost-minimising potential P1, whose failure is diagnosed
%%  rather than merely observed.  REFUTED: all five natural constructive rules
%%  for the witnessing transfer, under both the Psi and the weaker P3
%%  criterion -- in 2,264 of 7,416 crux partitions no transfer out of the
%%  envious agent's own bundle works at all.  So the repairing move is not
%%  local to the envy, and any proof must be non-constructive.  See also
%%  rem:descent-stronger: this route proves MORE than Conjecture conj:main.
%% ==========================================================================

\section{Approach 7: descent on a lexicographic potential}
\label{sec:approach7}

\subsection{The reformulation that suggests it}

Section~\ref{sec:approach5} characterises good allocations by the two-tier
condition: $\alloc$ is good precisely when there is a bipartition
$S \subseteq \agents$ with
\[
  \cost_i(\bundle{i}) - [\,i \in S\,]
  \;\le\;
  \cost_i(\bundle{j}) - [\,j \in S\,]
  \qquad \text{for all } i \ne j .
\]
Writing $\delta_j := [\,j \in S\,]$ and
$\hat\cost_i(\bundle{j}) := \cost_i(\bundle{j}) - \delta_j$, this says exactly
that \emph{every agent holds
a bundle minimising its discounted cost}, where the discount is $1$ on the paid
slots and $0$ elsewhere and depends only on the slot, not on the agent. Unpacked
by tier:

\begin{itemize}
  \item within $S$, and within $\agents \setminus S$: exact envy-freeness;
  \item an unpaid agent \emph{strictly} prefers its own bundle to every paid
    agent's bundle;
  \item a paid agent envies an unpaid agent by at most $1$.
\end{itemize}

Two consequences are worth stating before any search is set up. First, $S =
\agents$ and $S = \emptyset$ impose \emph{identical} constraints, since the two
indicators cancel; this is the structural reason the paid set is never all of
$\agents$, and it is the combinatorial shadow of
Proposition~\ref{prop:zero-coordinate}. Second, the existential over subsidy
vectors has become an existential over bipartitions, of which there are $2^n$ ---
the first formulation in which it is finite.

\subsection{The search space and the potential}

By Theorem~\ref{thm:hs-characterisation}(ii) an allocation is envy-freeable iff
it maximises welfare among reassignments of its own bundles. So we search over
\emph{partitions} and let the assignment be determined:

\begin{definition}[Canonical form]
\label{def:canonical}
For a partition $\mathcal B = (B_1,\dots,B_n)$ of $\items$, the \emph{canonical
allocation} $\alloc(\mathcal B)$ assigns bundles to agents by a minimum-cost
perfect matching. It is envy-freeable, so $\pathw{\alloc(\mathcal B)}$ is finite.
\end{definition}

\begin{definition}[The potential]
\label{def:potential}
For a partition $\mathcal B$ with canonical allocation $\alloc$, set
\[
  \Psi(\mathcal B) \;:=\;
  \Bigl(\; \max_{i \in \agents} \pathw{\alloc}(i), \;\;
           \sum_{i \in \agents} \abs{B_i}^2 \;\Bigr)
  \;\in\; \Z_{\ge 0} \times \Z_{\ge 0},
\]
ordered lexicographically. A \emph{transfer} moves one chore from one bundle to
another.
\end{definition}

The conjecture this approach rests on is a statement about single transfers.

\begin{conjecture}[Descent lemma]
\label{conj:descent}
Let $\mathcal B$ be a partition whose canonical allocation has
$\max_i \pathw{\alloc}(i) \ge 2$. Then some transfer $\mathcal B'$ satisfies
$\Psi(\mathcal B') < \Psi(\mathcal B)$.
\end{conjecture}

\subsection{The reduction, which is unconditional}

\begin{theorem}[Descent suffices]
\label{thm:descent-suffices}
Conjecture~\ref{conj:descent} implies Conjecture~\ref{conj:main} in full: every
instance with negative dichotomous valuations admits an envy-free solution
$(\alloc, \subsidy)$ with $\subsidy \in \set{0,1}^n$ and total subsidy at most
$n-1$, and one is computable in time $O(m^4 n^4)$ given value-oracle access.
\end{theorem}

\begin{proof}
Start from any partition and pass to canonical form. While
$\max_i \pathw{\alloc}(i) \ge 2$, apply Conjecture~\ref{conj:descent} to obtain a
transfer strictly decreasing $\Psi$.

\emph{Termination.} $\Psi$ takes values in $\Z_{\ge 0} \times \Z_{\ge 0}$ under
the lexicographic order, which is a well-order, so no infinite strictly
decreasing sequence exists. Quantitatively, every arc weight is at most $m$ in
absolute value and $\pathw{\alloc}$ is a weight of a simple path, so the first
coordinate lies in $\set{0,1,\dots,m}$; the second lies in $\set{m, \dots, m^2}$,
since $\sum_i \abs{B_i}^2$ is minimised by an equal split and maximised by
putting everything in one bundle. Hence $\Psi$ has at most $(m+1)(m^2+1) =
O(m^3)$ values and the loop runs $O(m^3)$ times.

\emph{Correctness.} The loop can only exit with $\max_i \pathw{\alloc}(i) \le 1$.
By Observation~\ref{obs:integrality} the minimum subsidy vector is integral and
non-negative, so $\optsubsidy \in \set{0,1}^n$; by
Theorem~\ref{thm:hs-minsubsidy} the pair $(\alloc, \optsubsidy)$ is an envy-free
solution; and by Proposition~\ref{prop:zero-coordinate} the total is at most
$n-1$.

\emph{Complexity.} One iteration examines $O(mn)$ transfers. Each requires a
minimum-cost perfect matching on $n$ agents, $O(n^3)$ by the Hungarian method
with $O(n^2)$ oracle calls, followed by an all-pairs longest-path computation on
an $n$-vertex graph, $O(n^3)$. So an iteration costs $O(mn \cdot n^3) = O(mn^4)$
and the whole procedure $O(m^4 n^4)$.
\end{proof}

Theorem~\ref{thm:descent-suffices} is the reason this approach is worth the
space: unlike every earlier route, the gap between what is proved and
Conjecture~\ref{conj:main} is a \emph{single} statement about \emph{single}
transfers, and closing it delivers the polynomial-time clause for free rather
than as a separate construction.

\subsection{Why this potential and not the obvious one}

The first potential tried was
$\Psi_1 = (\max_i \pathw{\alloc}(i), \#\set{i : \pathw{\alloc}(i) = \max},
\sum_i \cost_i(B_i))$, minimising total cost as a tie-break. It fails, and the
witness explains the failure rather than merely exhibiting it. At $n = 3$,
$m = 4$ with all four chores in one bundle and the other two empty,

\[
  \mathcal B = (\set{g_1,g_2,g_3,g_4}, \emptyset, \emptyset),
  \qquad \pathw{\alloc} = (0,0,2),
  \qquad \textstyle\sum_i \cost_i(B_i) = 2 ,
\]
every transfer leaves $\max_i \pathw{\alloc}(i) = 2$ and \emph{raises} the total
cost. The reason is that dichotomous costs cap: the loaded agent already values
its bundle at $2$, and removing one chore does not reduce that, while the
recipient's cost rises. So the third coordinate pins the search at a utilitarian
optimum --- precisely the class shown in Section~\ref{sec:approach6} not to be
good in general. The tie-break must reward \emph{spreading chores out}, not
minimising their total cost. Both $\sum_i \abs{B_i}^2$ and $\max_i \abs{B_i}$ do;
$\sum_i \pathw{\alloc}(i)$ in place of the maximum does not.

\begin{center}
\begin{tabular}{llr}
\toprule
potential & tie-break & stuck partitions \\
\midrule
$(\max \ell,\, \#\text{at max},\, \sum_i \cost_i(B_i))$ & total cost   & $449$ \\
$(\max \ell,\, \max_i \abs{B_i})$                       & largest bundle & $160$ \\
$(\sum_i \ell_i,\, \sum_i \abs{B_i}^2)$                 & total envy first & $464$ \\
\midrule
$(\max \ell,\, \#\text{at max},\, \max_i \abs{B_i})$    & ---          & $0$ \\
$(\max \ell,\, \#\text{at max},\, \sum_i \abs{B_i}^2)$  & ---          & $0$ \\
$\Psi = (\max \ell,\, \sum_i \abs{B_i}^2)$              & ---          & $0$ \\
\bottomrule
\end{tabular}
\end{center}

$\Psi$ is the shortest surviving potential, which is why
Definition~\ref{def:potential} takes it.

\subsection{Where the remaining burden sits}

The second coordinate of $\Psi$ moves in a completely explicit way, and this
splits Conjecture~\ref{conj:descent} into two unequal halves.

\begin{lemma}[Balance dichotomy]
\label{lem:balance-dichotomy}
Transferring a chore from a bundle of size $s$ to a bundle of size $t$ changes
$\sum_i \abs{B_i}^2$ by exactly $2(t - s + 1)$. Hence the transfer strictly
decreases the second coordinate of $\Psi$ if and only if $t \le s-2$.
Consequently, if $\mathcal B$ is \emph{balanced} --- all bundle sizes within $1$
of one another --- then no transfer decreases the second coordinate, and
Conjecture~\ref{conj:descent} demands a transfer that strictly decreases
$\max_i \pathw{\alloc}(i)$ itself.
\end{lemma}

\begin{proof}
The sizes change from $(s,t)$ to $(s-1,t+1)$, so the sum changes by
$(s-1)^2 + (t+1)^2 - s^2 - t^2 = 2(t-s+1)$, which is negative exactly when
$t < s-1$. If all sizes lie in $\set{q, q+1}$ then $t \ge q \ge s-1$ for every
ordered pair, so no transfer qualifies.
\end{proof}

So Conjecture~\ref{conj:descent} is really two claims. On unbalanced partitions
one must find a size-equalising transfer that does not \emph{increase} the
maximum subsidy; on balanced partitions one must strictly decrease the maximum
subsidy by a single transfer. The second is the hard half and the natural next
target: it is a statement about balanced partitions only, where every bundle has
size $\lfloor m/n \rfloor$ or $\lceil m/n \rceil$, which is a far more rigid object than an
arbitrary allocation. Note that the transfer is followed by re-matching, so a
single transfer may permute the assignment substantially --- this is where the
strength of the move comes from, and any proof must use it.

\subsection{The obstruction: the witnessing transfer is not local}
\label{sec:approach7-obstruction}

Conjecture~\ref{conj:descent} asserts that \emph{some} transfer works. A proof
wants a \emph{rule} --- a named transfer, definable from the envy structure,
that always works. The obvious candidates all fail, and the way they fail is
informative.

Two mechanisms can reduce an arc $\edgew{\alloc}(i,j) = \cost_i(\bundle{i}) -
\cost_i(\bundle{j})$: remove a chore from $\bundle{i}$ to lower
$\cost_i(\bundle{i})$, or add one to $\bundle{j}$ to raise
$\cost_i(\bundle{j})$. Under dichotomous costs \emph{both can fail outright},
because costs cap and marginals need not be monotone:

\begin{itemize}
  \item $\cost(S) = \min(\abs{S},1)$ has no chore whose removal lowers the cost
    of a two-element bundle;
  \item $\cost(S) = \max(\abs{S}-1,0)$ has no chore whose addition raises the
    cost of the empty bundle.
\end{itemize}

Note that the second example also rules out the submodular shortcut: a chore may
have marginal $1$ over a superset of $\bundle{j}$ and marginal $0$ over
$\bundle{j}$ itself, so ``some chore raises $\cost_i(\bundle{j})$'' cannot be
inferred from $\cost_i$ being large somewhere. This is the same failure of the
exchange property that has obstructed every earlier approach in this report.

Five rules were tested against all $7{,}416$ crux partitions --- balanced, with
$\max_i \pathw{\alloc}(i) \ge 2$ --- under two criteria: the $\Psi$ criterion
(strictly reduce $\max_i \pathw{\alloc}(i)$) and the weaker $P_3$ criterion
(reduce $(\max_i \pathw{\alloc}(i), \#\text{at max})$ lexicographically).

\begin{center}
\begin{tabular}{llrr}
\toprule
rule & transfer & $\Psi$ criterion & $P_3$ criterion \\
\midrule
R1 & argmax-$\ell$ bundle $\to$ max-path endpoint & $2{,}856$ & $3{,}168$ \\
R2 & argmax-$\ell$ bundle $\to$ its cheapest bundle & $3{,}827$ & $4{,}021$ \\
R3 & argmax-$\ell$ bundle $\to$ anywhere           & $5{,}112$ & $5{,}152$ \\
R4 & anywhere $\to$ max-path endpoint              & $6{,}039$ & $6{,}467$ \\
R5 & largest bundle $\to$ smallest bundle          & $3{,}972$ & $4{,}086$ \\
\midrule
   & \emph{every transfer} (Conjecture~\ref{conj:descent}) & $7{,}416$ & $7{,}416$ \\
\bottomrule
\end{tabular}
\end{center}

The decisive row is R3. In $2{,}264$ of $7{,}416$ crux partitions, no transfer
\emph{out of the envious agent's own bundle} works, even under the weaker
criterion --- the witnessing move involves two bundles neither of which belongs
to an agent attaining the maximum. So the move that repairs the envy is not
local to the envy. Any proof of Conjecture~\ref{conj:descent} must therefore be
non-constructive, or must name a rule of a kind not yet imagined; the five
natural ones are exhausted.

\begin{remark}[The descent lemma is strictly stronger than what is needed]
\label{rem:descent-stronger}
Conjecture~\ref{conj:descent} asserts more than
Conjecture~\ref{conj:main}: it says both that good allocations exist and that
greedy descent on $\Psi$ finds one, i.e.\ that $\Psi$ has no local minimum with
$\max_i \pathw{\alloc}(i) \ge 2$. Conjecture~\ref{conj:main} is the statement
that the \emph{global} minimum of $\Psi$ has first coordinate at most $1$. The
extra strength is what buys the polynomial-time clause in
Theorem~\ref{thm:descent-suffices} at no cost, but it also means this route is
harder than the target, and the obstruction above may be a symptom of exactly
that. A proof of Conjecture~\ref{conj:main} alone need not go through any local
search.
\end{remark}

\subsection{Evidence}

Conjecture~\ref{conj:descent} was tested by enumerating \emph{every} partition of
every test instance, computing $\Psi$ at each and at all $O(mn)$ of its
transfers, and recording any partition with $\max_i \pathw{\alloc}(i) \ge 2$
admitting no strict decrease.

\begin{center}
\begin{tabular}{lrr}
\toprule
family & instances & stuck partitions \\
\midrule
random, $n \le 4$, $m \le 6$                    & $1{,}732$ & $0$ \\
EF-free only, $n \le 5$, $m \le 6$              & $253$     & $0$ \\
R10 Set-Splitting, certified hard               & $4$       & $0$ \\
random, $n = 3$, $m \le 9$                      & $798$     & $0$ \\
random, $n = 4$, $m \le 7$                      & $280$     & $0$ \\
random, $n = 5$, $m \le 6$                      & $75$      & $0$ \\
\midrule
\emph{crux partitions only} (balanced, $\max \ell \ge 2$) & $7{,}416$ & $0$ \\
\bottomrule
\end{tabular}
\end{center}

The last row is the one that matters, and it was measured separately because the
others do not imply it: ``zero stuck partitions over all partitions'' would be
consistent with the crux population being empty. It is not. It is empty when
$m = n$ --- a balanced partition then gives every agent one chore, and the
min-cost matching caps $\pathw{\alloc}$ at $1$, so there is nothing to prove ---
and it grows with $m$, reaching $990$ crux partitions at $n=3$, $m=9$ and
$2{,}112$ at $n=4$, $m=6$. Every one of them admits a $\Psi$-decreasing transfer.

The large-$m$ rows matter specifically: a size-based potential should fail when
bundles grow, since $\sum_i \abs{B_i}^2$ can only be decreased by equalising
sizes while a large bundle may be cheap under a capped cost function. No such
failure appears up to $m = 9$. Separately, no instance in any of these families
lacked a good allocation, so Conjecture~\ref{conj:main} itself remains unrefuted
throughout.
